% Folder 134-2, batch 07 — pages 121–140 (position in the folder; see below).
% Pass: Opus 5.5 (claude-opus-5-5), 2026-10-15 — first pass, unchecked against the pages by a human.
% Revised: Opus 5.5 (claude-opus-5-5), 2026-10-15 — the grey pencil of p. 122 identified as R. Brown's by Michel Hua on the facsimile (#35).
%
% Special rules for folder 134-2 (issues #44 and #57, the latter as decided by the
% user on 2026-10-15) applied: Pursuing Stacks text is not transcribed; his own
% unprinted letters are transcribed in full; texts sent to him are transcribed in
% full; texts by others not addressed to him are summarised, one \page{} and one
% \note{} per page, any mark in Grothendieck's hand given in full.
%
% Numbering: \page{} numbers the position in the folder, as in batch 6 (and as the
% tiles and the facsimile do). From p. 115 on the archivists' pencil runs one ahead
% of the position (no sheet pencilled 115 is in the scan): here p. 121 is pencilled
% 122, p. 140 is pencilled 141.
%
% Contents of the batch:
%  - p. 121 (pencilled 122): blank; absent.
%  - pp. 122-126 (pencilled 123-127): R. Brown's typed English translation,
%    « Translation of a letter of A. Grothendieck to L. Breen 17/2/1975 », 5 pp.,
%    complete, signed « Alexandre. ». Others' text, not addressed to him:
%    summarised. The French typed copy of this letter is transcribed in batch 3
%    (pp. 48-55). Marks: « topoi » and « stack » in grey pencil in the left margin
%    of p. 122, R. Brown's grey pencil (checked on the facsimile by Michel Hua,
%    2026-10-15); given in the note, not attributed to Grothendieck.
%    « this seems suspect » in the margin of p. 123 is typed (the translator's).
%  - p. 127 (pencilled 128): unrelated verso (computer listing); absent.
%  - pp. 128-140 (pencilled 129-141): R. Brown's typed translation « Translation
%    of an extract from a Letter from A. Grothendieck to L. Breen 17-19 July,
%    1975 », his pp. 1-13; runs on past p. 140. Summarised, one note per page.
%    Two words of the title are blanked out with correction fluid. The French
%    handwritten original of this letter (pp. 145-184) falls in later batches.
%    No mark in Grothendieck's hand found on these pages.
%
% Letters in this batch: none of his in his own words — only Brown's translations
% of the letters to Breen of 17.2.1975 and 17-19.7.1975 (summarised).
\documentclass[11pt,a4paper]{article}
\input{../preamble/grothendieck}

\folder{134-2}
\batch{7}
\pages{121}{140}
\foldertitle{[Chapitre I : Take off (pages 1 à 65) et table des matières provisoire] : tapuscrits annotés (19/02-22/02/1983), note manuscrite (s.d.), copies de lettre (1975, s.d.).}
\dating{1975-[1983]}
\watermark{Édition de démonstration}

\begin{document}

\section{Traduction anglaise (R. Brown) de la lettre à L. Breen du 17.2.1975}

\page{122}
\note{numérotation : à partir de la p.~115, le crayon des archivistes a un de plus que la position dans le dossier ; cette page est crayonnée 123. Début de la traduction anglaise dactylographiée, par R. Brown, de la lettre à L. Breen du 17.2.1975 (« Translation of a letter of A. Grothendieck to L. Breen 17/2/1975 », « Dear Larry »), dont la copie française est transcrite au lot~3 (pp.~48 à 55) ; texte d'un tiers, résumé. Contenu : le pro-type d'homotopie associé à un topos ; critère d'Artin--Mazur pour qu'un morphisme de topos soit une « équivalence d'homotopie en dim $\leqslant n$ », reformulé en termes de $n$-champs localement constants (énoncés (A) et (B)) ; énoncé (C) : les $(n-1)$-champs localement constants sur $X$ comme $n$-foncteurs du $n$-groupoïde fondamental $\Pi_n(X)$ vers $((n-1)\text{-Cat})$ ; cas $n = 1$ (Poincaré, SGA~3). En marge gauche, au crayon gris, de la main de R. Brown (vérifié sur le fac-similé par Michel Hua, 2026-10-15) : « topoi » (en face de « toposes ») et « stack », souligné (en face de « n-champ »).}

\page{123}
\note{p.~« 2 » de la traduction. Le cas $n=1$ : le groupoïde $C_1$ retrouvé à partir du topos « multigalois » des systèmes locaux ; le $n$-topos classifiant $\mathbb{C}_n = (n\text{-}\mathit{Hom})(C_n, ((n-1)\text{-Cat}))$ (« Brrr! ») ; doute sur l'existence d'un topos classifiant ordinaire pour un $n$-groupoïde dès $n = 2$, rapproché de l'échec présumé du théorème de van Kampen pour $\Pi_2$. En marge, dactylographié et coupé à gauche : « [t]his seems / suspect » (remarque du traducteur, non de Grothendieck).}

\page{124}
\note{p.~« 3 » de la traduction. Cohomologie d'un $n$-groupoïde $C_n$ via (B) : $\Gamma_X(F) \simeq \mathit{Hom}(e^X_{n-1}, F)$ ; « apologie » à la lettre précédente : cohomologie non commutative d'un topos par intégration de $n$-champs ; pour un $n$-champ de Picard strict défini par un complexe $L^\bullet$, $H^i(X, L^\bullet) = \pi_{n-i}\Gamma_X(F)$, et le foncteur « espace classifiant » $\mathbb{C}$ pour les $H^i$ de degré quelconque ; début d'un « N.B. » sur le morphisme canonique $\mathbb{C}(\Gamma_X F) \to \Gamma_X(\mathbb{C}F)$.}

\page{125}
\note{p.~« 4 » de la traduction. Fin du N.B. ; second N.B. : pour un $n$-champ en groupoïdes, $H^i(X,F) = \pi_{n-i}(\Gamma_X(F))$ sans structure de Picard, et pour un $n$-Gr-champ jusqu'à $H^{n+1}$ par le $(n+1)$-champ classifiant ; cas d'un $n$-champ localement constant donné par un $(n+1)$-foncteur $C_{n+1} \to$ (strict Picard $n$-Cat), définition de $H^i(C, F)$ pour un $\infty$-groupoïde $C$, et question : décrire la flèche $H^{n+1}(X,F) \simeq H^{n+1}(\Pi_{n+1}(X), F)$.}

\page{126}
\note{p.~« 5 » de la traduction, fin de la lettre, signée « Alexandre. ». Énoncés (A') et (B') pour un $(n+1)$-foncteur $f_{n+1}\colon C_{n+1} \to D_{n+1}$ ; la cohomologie d'un topos par intégration de champs, sans complexes de faisceaux ni résolutions injectives ; les trois points de vue (a) complexes de faisceaux (algèbre homologique commutative), (b) Čechiste ou semi-simplicial (algèbre homotopique), (c) $n$-champs (algèbre homologique non commutative) — en (a) on résout les coefficients, en (b) la base, en (c) ni l'un ni l'autre.}

\section{Traduction anglaise (R. Brown) d'un extrait de la lettre à L. Breen des 17--19.7.1975}

\page{128}
\note{début de la traduction anglaise dactylographiée, par R. Brown, d'un extrait de la lettre à L. Breen des 17--19 juillet 1975 (« Translation of an extract from a Letter from A. Grothendieck to L. Breen 17-19 July, 1975 ») ; deux mots du titre sont effacés au correcteur. Texte d'un tiers, résumé ; l'original français manuscrit est plus loin dans le dossier (pp.~145 à 184). Contenu : réponse à la construction de Breen des $n$-catégories non strictes et du foncteur nerf (mérite d'exister, mais technique et sans foncteur des ensembles (semi-)simpliciaux vers les $n$-groupoïdes) ; idée d'une « liste de cadeaux de Noël » menant à une description axiomatique ; l'idée de Segal non comprise ; ce qu'il demande aux $n$-catégories de Picard strictes.}

\page{129}
\note{p.~« 2 » de la traduction. Les $n$-catégories de Picard strictes équivalentes aux complexes de chaînes de longueur $n$ ; la « grosse bévue » sur les H-espaces rectifiée par Breen ; question d'un analogue du « tapis » de Dold--Puppe pour les groupes semi-simpliciaux non commutatifs, et le yoga « simplicial sets $\longleftrightarrow$ $\infty$-groupoids » comme version ensembliste ultime de Dold--Puppe ; les groupes dans (1-cat), Gr-catégories strictes (livre de Mme Sinh), calcul de fractions.}

\page{130}
\note{p.~« 3 » de la traduction. Les groupes dans (1-cat) décrits par un quadruplet $(L_1, L_0, d, \theta)$, $d\colon L_1 \to L_0$, $\theta\colon L_0 \to \mathrm{Aut}_{\mathrm{Gr}}(L_1)$, avec les deux axiomes (a) $d(\theta(x_0).x_1) = \mathrm{int}(x_0)\, d(x_1)$, (b) $\theta(d(x_1)) = \mathrm{int}(x_1)$ (modules croisés) ; $\pi_0 = \operatorname{Coker} d$, $\pi_1 = \operatorname{Ker} d$ central, invariant de Postnikov--Sinh $\alpha \in H^3(\pi_0, \pi_1)$ ; rencontres de ces objets, notamment pour la classification des groupes formels sur un corps parfait.}

\page{131}
\note{p.~« 4 » de la traduction. Question de la généralisation aux groupes dans ($n$-cat) : un complexe de chaînes non commutatif $L_n \to L_{n-1} \to \dots \to L_1 \to L_0 \to 1$ avec des structures supplémentaires du type de $\theta$, modélisant le type d'homotopie tronqué des groupes topologiques ; « Have you candidates to propose? ». Puis, entre crochets, du traducteur : un commentaire sur la bidualité et l'homologie est omis. Bas de page blanc.}

\page{132}
\note{p.~« 5 » de la traduction. Sur la question 7 de Breen : ce qui reste d'un topos $X$ dans la catégorie homotopique — les $\pi_i$, la cohomologie à coefficients constants, et précisément les $n$-champs localement constants ; triangle d'objets qui se déterminent mutuellement : topos (ou espace, ou complexe semi-simplicial) modulo $n$-homotopie, $n$-groupoïdes, $n$-catégories « spéciales » (à $n$-équivalence près), dessiné à la machine avec trois doubles flèches.}

\page{133}
\note{p.~« 6 » de la traduction. Définition des $n$-catégories « spéciales » (ou $n$-galoisiennes), groupoïde fondamental $G_n$, systèmes locaux, réalisation géométrique ; $G_n$ comme sous-$n$-catégorie de $\mathrm{Hom}_n(E_n, (n\text{-cat}))$ formée des foncteurs exacts ; hypothèse « localement homotopiquement trivial » ; pour le topos étale d'un schéma, pro-$n$-groupoïde fondamental et ind-$n$-catégorie.}

\page{134}
\note{p.~« 7 » de la traduction. La « résistance instinctive » de Breen à réduire un topos à son type d'homotopie ; vers une notion plus fine de type d'homotopie : faisceaux constructibles sur le topos étale d'un schéma, partitions finies en strates localement fermées, faisceaux « $P$-constructibles ».}

\page{135}
\note{p.~« 8 » de la traduction. Stabilité par les opérations ($R\,\mathrm{Hom}$ (souligné), $Rj_*j^*$) et stratifications équisingulières ; exemple de deux strates $X_0$ fermée et $X_1 = X \setminus X_0$ : un faisceau $F$ donné par $F_1$ sur $X_1$ et un morphisme $F_0 \to i_0^* i_{1*}(F_1)$, exprimable par les seuls types d'homotopie de $X_0$ et $X_1$.}

\page{136}
\note{p.~« 9 » de la traduction. Le foncteur exact à gauche $\phi$ ; « voisinage tubulaire étale » de $X_0$ ; sur $\mathbb{C}$, voisinage tubulaire $T$ de $X_0$ dans $X$, bord $\partial T$, $\mathring{T} = T \setminus X_0 \simeq \partial T \times [0,1]$, $\widetilde{X}_1 = X_1 \setminus \mathring{T}$, et $X$ reconstitué à homéomorphisme près à partir du diagramme $X_0 \xleftarrow{p} \partial T \xrightarrow{j} \widetilde{X}_1$ ($p$ fibration, $j$ inclusion).}

\page{137}
\note{p.~« 10 » de la traduction. $X$ comme somme amalgamée ; $\phi(F_1) \simeq p_* j^*(\widetilde{F}_1)$ ; donnée de $F$ équivalente à $(F_0, \widetilde{F}_1, \tilde{u}\colon p^*(F_0) \to j^*(\widetilde{F}_1))$ ; variantes du diagramme avec $\mathring{T} \xrightarrow{j'} X_1$, $p'\colon \mathring{T} \to X_0$, puis $p''\colon \mathring{T} \to T$.}

\page{138}
\note{p.~« 11 » de la traduction. $p''$ « moralement » une fibration à jolie fibre ; les faisceaux $P$-constructibles pour une stratification équisingulière ramenés à un diagramme de types d'homotopie ordinaires $X_0 \xleftarrow{p} \Delta \xrightarrow{j} X_1$ et des systèmes de coefficients locaux reliés par $p^*(F_0) \to j^*(F_1)$ ; exercice proposé : exprimer les opérations usuelles dans ce dictionnaire.}

\page{139}
\note{p.~« 12 » de la traduction. Diagrammes vus dans la catégorie homotopique des ensembles semi-simpliciaux ; la « topologie modérée » (\emph{tame topology}) : une stratification équisingulière d'un espace modéré compact donne canoniquement un diagramme de variétés à bord, reliées par des fibrations localement triviales et des inclusions de bord, d'où une description par dévissage. Un blanc dans la frappe avant « the inclusions of these pieces ».}

\page{140}
\note{p.~« 13 » de la traduction. Les faisceaux $P$-constructibles exprimés sur ce diagramme ; les six opérations traduites dans ce cadre homotopique ; pour un morphisme modéré $f\colon X \to Y$, les quatre opérations $Rf_*$, $Rf_!$, $Lg^*$, $g^!$ ; paraphrase attendue pour les schémas excellents ; le « type d'homotopie fin » comme limite des « $P$-types d'homotopie ». La phrase continue au-delà du lot.}

\end{document}
